\section{模型假设}

表\ref{tab:assumptions}只列求解所必需、而题目未直接给定的边界。每项假设都对应程序口径和失效条件；新证据若越过这些条件，必须重算相关模型，不能只修改文字解释。

{\small
\setlength{\tabcolsep}{4pt}
\renewcommand{\arraystretch}{1.08}
\begin{longtable}{p{0.15\textwidth}>{\raggedright\arraybackslash}p{0.49\textwidth}>{\raggedright\arraybackslash}p{0.25\textwidth}}
\caption{模型假设、固定口径与失效处理}\label{tab:assumptions}\\
\toprule
对象 & 固定口径 & 失效时的处理 \\
\midrule
\endfirsthead
\multicolumn{3}{c}{续表\thetable}\\
\toprule
对象 & 固定口径 & 失效时的处理 \\
\midrule
\endhead
\midrule
\multicolumn{3}{r}{续下页}\\
\endfoot
\bottomrule
\endlastfoot
预测信息集 & $560/140$ 条记录构成时间外边界；测试标签不参与训练、调参、早停或缺失值统计。同场赛后字段只可进入严格滞后的球队历史，同日比赛在日末统一更新。内部以百万人拟合，正式赛程需求再乘 $10^6$ 转为人。 & 字段获得时间晚于开球、验证窗标签进入滞后量或单位不一致时，重新构造特征并完成扩窗验证。 \\
\addlinespace
联合排程域 & 比赛、场馆和时段是一个联合决策。当地日期由 UTC 转换；第二、三轮旅行以上一轮全部安保合格候选场馆为起点取均值，遵守题定口径。归一化边界来自同一硬可行域；未证最优的辅助问题采用 incumbent 与 best bound 的保守包络。 & 只有题方更改旅行定义、候选场馆或硬约束时才重建域；不得把现实中未提供的禁用日或合同条件暗设为硬约束。 \\
\addlinespace
第三轮状态 & $24$ 场共享第二轮结束后的同一状态截面。每个样本联合生成全部比分并统一排名，跨线同分块按比例分配，使晋级名额质量恒为 $32$；更新后的泊松均值截断在 $[0.15,4.50]$。 & 若状态边界长期饱和，应重估状态函数；增加模拟次数不能修复参数错配。逐场滚动更新只能作为另一个情景。 \\
\addlinespace
静态与动态方案 & 静态动作按赛前信息优化后锁定；动态动作只在更新环境中重选资源。两者使用共同动作域、价格边界、归一化边界和日预算，因此动态最优值不得低于锁定静态值。$-1<\varepsilon_i<0$ 时票价取解析右端点。 & 若出现 $\varepsilon_i\leq-1$，票价改用有界一维搜索；同界下动态值更低则说明实现或接口错误。 \\
\addlinespace
现实比较 & $2022$ 赛程取自保存的 OpenFootball 记录并按 FIFA 页面核验，球队实力使用开赛前完整 Elo 快照。小时、公里、场馆日等结构量直接复算；商业、成本和安保缺失量仅在题内最近邻代理层使用。 & 代理集中或邻居替换只降低综合分的证据强度，不能改写结构事实；补齐真实运营数据后须重做代理层。 \\
\addlinespace
随机性与复现 & 主种子为 $20260813$，晋级模拟一次生成 $20000\times24$ 样本；CP-SAT 单线程运行，正式 CSV 再由独立程序复算硬约束。 & 样本名额不守恒、硬约束失败或不同环境输出不一致时，停止解释结果并回查数据与求解接口。 \\
\end{longtable}
}

这些假设不把未知事实伪装成观测。题定口径用于主模型，现实缺失量保留代理标记；因此论文中的结论强度始终受数据来源和可复算范围约束。
